
AIS automaatne identifitseerimise süsteem (\emph{ingl. k} \emph{Automatic Identification System)}

ARM Arm Holdings'i poolt arendatav vähendatud käsustikuga arvutiarhitektuur (\emph{ingl. k Advanced RISC Machines)}

CNN konvolutsiooniline närvivõrk (\emph{ingl. k} \emph{convolutional neural network)}

COCO Tuntud objektid kontekstis, annoteeritud fotode andmestik (\emph{ingl. k} \emph{Common}

\emph{Objects in Context)}

CPU Arvuti protsessor \emph{(ingl. k} \emph{Central Processing Unit)}

CSI Kaamera jadaühenduse liides (\emph{ingl. k} \emph{Camera Serial Interface)}

CUDA Arvutamise ühtne seadmearhitektuur (\emph{ingl. k} \emph{Compute Unified Device}

\emph{Architecture)}

DL Süvaõpe (\emph{ingl. k} \emph{Deep Learning)}

FPS Kaadrid sekundis (\emph{ingl. k} \emph{Frames Per Second)}

FSD Täielikult autonoomne süsteem Tesla sõidukitele (\emph{ingl. k} \emph{Full Self Driving)}

GB gigabait (\emph{ingl. k} \emph{GigaByte)}

GFLOPs miljardit ujukomarvu operatsiooni sekundis (\emph{ingl. k} \emph{Billions of Floating point}

\emph{Operations per second)}

GPU graafikatöötlusseade, graafikakaart (\emph{ingl. k} \emph{Graphics Processing Unit)}

Gbps gigabitti sekundis (\emph{ingl. k} \emph{Gigabits per second)}

HAVS käsivarre vibratsiooni sündroom (\emph{ingl. k} \emph{Hand Arm Vibration Syndrome)}

HMMA 16 komakohalise ujukomaarv maatriksi korrutamine ja summeerimine (\emph{ingl. k}

\emph{Half precision Matrix Multiply and Accumulate)}

IMMA täisarvulise maatriksi korrutamine ja summeerimine (\emph{ingl. k} \emph{Integer Matrix}

\emph{Multiply and Accumulate)}

L4T Linuxi operatsioonisüsteem Nvidia Tegra manussüsteemidega seadmetele (\emph{ingl.}

\emph{k} \emph{Linux for Tegra)}

MB megabait (\emph{ingl. k} \emph{megabyte)}

NVDEC Nvidia videodekooder (\emph{ingl. k} \emph{Nvidia Video Decoder)}

NVDLA~Nvidia süvaõppe kiirendi (\emph{ingl. k} \emph{Nvidia Deep Learning Accelerator)}

ONNX Avatud värvivõrkude vahetamise formaat (\emph{ingl. k} \emph{Open Neural Network}

\emph{Exchange)}

POE toide üle võrgukaabli (\emph{ingl. k} \emph{Power Over Ethernet)}

PVA Programmeeritav Masinnägemise Kiirendi (\emph{ingl. k} \emph{Programmable Vision}

\emph{Accelerator)}

SBC monoplaatarvuti (\emph{ingl. k} \emph{single board computer)}

SoM süsteem moodulil või süsteem kiibistikul (\emph{ingl. k} \emph{System on Module or SOC}

\emph{system on chip)}

TFLOPs triljonit ujukomaarvu operatsiooni sekundis (\emph{ingl. k} \emph{Trillions of Floating point}

\emph{Operations per second)}

TOPs triljonit operatsiooni sekundis (\emph{ingl. k} \emph{Trillions of Operations per second)}

USB universaalne jadasiin (\emph{ingl. k} \emph{Universal Serial Bus)}

YOLO sa vaatad ainult korra, filosoofia, kus närvivõrk töötleb pilti ainult ühe korra

~(\emph{ingl. k} \emph{You Only Look Once)}

eMMC integreeritud multimeediumi kaart (\emph{ingl. k} \emph{embedded MultiMedia Card)}

mAP täpsuste keskväärtuste keskmine (\emph{ingl. k} \emph{mean Average Precision)}

ms millisekund (\emph{ingl. k} \emph{millisecond)}

px piksel (\emph{ingl. k pixel})
